% Generated from locale/tr/content/proof-theory/cut-elimination/midsequent.tex; do not hand-edit.


\olfileid[tr]{pt}{cut}{mid}

\olsection{Orta Sekant Teoremi ve Herbrand Teoremi}

Kesme-giderimi teoreminin önemli bir sonucu orta sekant teoremidir. Bu teorem,
bir $\Gamma \Sequent \Delta$ sekantının !!a{proof}{}ı~$\pi$ varsa ve bu
sekanttaki her !!{formula} önek biçimindeyse, $S$'nin üstündeki bütün çıkarımlar
önermesel ve $S$'nin altındaki bütün çıkarımlar niceleyici çıkarımları olacak
biçimde, (\emph{orta sekant} denen) $S \ident \Pi \Sequent \Lambda$ sekantını
içeren bir $\pi'$ !!a{proof}{}ı bulunduğunu söyler. (Bir $!A$ !!^a{formula}{}ü
\[
\mathsf{Q}_1x_1\dots\mathsf{Q}_nx_n\,!B(x_1,\dots,x_n)
\]
biçimindeyse ve her $\mathsf{Q}_i$, $\lforall$ ya da~$\lexists$ ise
\emph{önekli} (ya da \emph{önek biçiminde})dir.) Orta sekant teoreminin
önemli bir sonucu da Herbrand teoremidir. Genel durumu betimlemek biraz
güçtür; fakat çok genel olarak şu biçimi alır: birinci dereceden mantığın her
!!{provable} $!A$ !!{sentence}{}si için, $!A$'yı !!{prove}{}maya yarayan
bir $!A'$ önermesel totolojisi vardır. $!B$'nin niceleyicisiz
olduğu $\lexists[x][!B(x)]$ biçimindeki !!{sentence}s için sürümü şöyledir:

\begin{thm}[Herbrand Teoremi]
$!B(x)$'in niceleyicisiz ve $\Proves \lexists[x][!B(x)]$ olduğunu varsayalım.
O zaman $\Proves !B(t_1) \lor \dots \lor !B(t_n)$ olacak biçimde $t_1$,
\dots, $t_n$ terimleri vardır.
\end{thm}

$!B(t_1) \lor \dots \lor !B(t_n)$ !!{formula}{}üne
$\lexists[x][!B(x)]$'in bir \emph{Herbrand ayrışımı} denir. $!B$
niceleyicisiz olduğundan bu ayrışım da niceleyicisizdir. Dolayısıyla
!!{provable} ise kesmesiz bir ispatla !!{provable}{}dir ve bunun sonucu olarak
yalnızca önermesel çıkarımlar kullanan !!a{proof}{}ı vardır. Öyleyse bir
totolojidir.

Herbrand Teoremi'nin güç kısmı, her !!{provable}
$\lexists[x][!B(x)]$ !!{formula}{}ü için bir Herbrand ayrışımının var olduğunu
göstermektir. Tersi, yani $\lexists[x][!B(x)]$ bir Herbrand ayrışımına sahipse
ispatlanabilir olduğu kolaydır. Gerçekten Herbrand ayrışımı
$\lexists[x][!B(x)]$'i gerektirir. Örneğin $n=2$ durumu için~\Log{G3c}'de bir
!!a{proof} şöyledir:
\[
\Axiom$!B(t_1) \fCenter \lexists[x][!B(x)], !B(t_1), !B(t_2)$
\Axiom$!B(t_2) \fCenter \lexists[x][!B(x)], !B(t_1), !B(t_2)$
\RightLabel{\LeftR{\lor}}
\BinaryInf$!B(t_1) \lor !B(t_2) \fCenter \lexists[x][!B(x)], !B(t_1), !B(t_2)$
\RightLabel{\RightR{\lexists}}
\UnaryInf$!B(t_1) \lor !B(t_2) \fCenter \lexists[x][!B(x)], !B(t_2)$
\RightLabel{\RightR{\lexists}}
\UnaryInf$!B(t_1) \lor !B(t_2) \fCenter \lexists[x][!B(x)]$
\DisplayProof
\]

$\Sequent \lexists[x][!B(x)]$'in bir $\pi$ !!a{proof}{}ından Herbrand
ayrışımını çıkarmak için, kesmesiz bir $\pi$ !!{proof}{}ındaki bütün
$\RightR{\lexists}$ çıkarımlarının sonda bulunduğu durumu ele alın:
\[
\AxiomC{}
\RightLabel{$\pi_1$}
\Deduce$ \fCenter \lexists[x][!B(x)], !B(t_1), \dots, !B(t_n)$
\RightLabel{\RightR{\lexists}}
\UnaryInf$ \fCenter \lexists[x][!B(x)], !B(t_2), \dots, !B(t_n)$
\RightLabel{$(n-1) \times \RightR{\lexists}$}
\Deduce$ \fCenter \lexists[x][!B(x)]$
\DisplayProof
\]
Bunlar $\pi$ içinde $\lexists[x][!B(x)]$'in etkin olduğu bütün
$\RightR{\lexists}$ çıkarımlarıysa $\pi_1$ içinde başka niceleyici çıkarımı
yoktur. Çünkü $!B(x)$ başka niceleyici içermez ve son sekantın solunda hiçbir
!!{formula} bulunmaz. Başka bir deyişle $ \fCenter \lexists[x][!B(x)],
!B(t_1), \dots, !B(t_n)$, $\pi$'nin orta sekantıdır. Dahası, orta sekantın
üstünde niceleyici çıkarımı bulunmadığından onun üstündeki bütün sekantlar
$\lexists[x][!B(x)]$'i içerir ve bunlar ne etkin ne de ana formüldür. Böylece
$\pi_1$ !!{proof}{}ından silinebilirler; $\fCenter !B(t_1), \dots, !B(t_n)$'nin
bir !!a{proof}{}ını ve $\RightR{\lor}$'un $n-1$ uygulamasıyla Herbrand
ayrışımı~$\Sequent !B(t_1) \lor \dots \lor !B(t_n)$'nin bir !!a{proof}{}ını
elde ederiz.

Sorun şudur: $\Sequent \lexists[x][!B(x)]$'in kesmesiz bir !!{proof}{}ında bile
bütün \RightR{\lexists} çıkarımlarının sonda tek bir blok oluşturduğunu
varsayamayız. Bu \RightR{\lexists} çıkarımları arasında bazı $!B(t_i)$'lerin
ana formül olduğu çıkarımlar bulunabilir. Orta sekant teoremine bu nedenle
ihtiyaç duyarız.

\begin{thm}[Orta Sekant Teoremi]\ollabel{thm:midsequent} Bir $\Gamma
  \Sequent \Delta$ sekantının \Log{G3c}-!!{proof}{}ı~$\pi$ varsa ve $\Gamma$
  ile $\Delta$ içindeki bütün !!{formula}s önek biçimindeyse,
  altında kalan bütün çıkarımlar niceleyici çıkarımları ve üstünde kalan bütün
  çıkarımlar önermesel çıkarımlar olacak biçimde bir $\Pi \Sequent \Lambda$
  sekantını içeren bir $\pi'$ !!a{proof}{}ı vardır.
\end{thm}

\begin{proof}
  $\pi$ içindeki bir niceleyici çıkarımının \emph{derecesini}, altında kalan
  önermesel çıkarımların sayısı; $\pi$'nin derecesi~$o(\pi)$'yi de bütün
  niceleyici çıkarımlarının derecelerinin toplamı olarak tanımlayın.
  $o(\pi) = 0$ ise hiçbir niceleyici çıkarımının altında önermesel bir çıkarım
  yoktur ve işimiz bitmiştir: Bir niceleyici çıkarımı varsa, bütün niceleyici
  çıkarımlarının tek öncülü olduğundan en üstte tek bir niceleyici çıkarımı
  vardır; onun öncülü orta sekanttır. Niceleyici çıkarımı yoksa son sekantın
  kendisini orta sekant olarak alırız.

  $o(\pi) > 0$ ise derecesi $>0$ olan en alttaki niceleyici çıkarımlarından
  birini seçin. Derecesi $>0$ olduğundan altında bir önermesel çıkarım
  bulunmalıdır. Seçilen çıkarım en alttaki olduğundan sonucu bir niceleyici
  çıkarımının öncülü olamaz: aksi hâlde ikinci, daha alttaki niceleyici
  çıkarımının da altında bir önermesel çıkarım bulunurdu. Hangi niceleyici
  çıkarımına ve onun altında hangi önermesel çıkarıma sahip olduğumuza göre
  (pek çok!) durumu ayırırız. Son sekant yalnızca önek biçimindeki formüllerden
  oluştuğundan niceleyici çıkarımının ana formülü, onu izleyen önermesel
  çıkarımda etkin olamaz. Aksi hâlde o çıkarımın ana formülü, bir önermesel
  bağlacın kapsamında bir niceleyici içerirdi. Alt-!!{formula} özelliği gereği
  bunun, son sekanttaki !!{formula}s{}den !!a{formula}{}ün alt-formülü olması
  gerekirdi.
  Öyleyse niceleyici çıkarımının ana !!{formula}{}ü, daha alttaki önermesel
  çıkarımın bağlamının bir !!a{element}{}ıdır.

  İşte iki örnek:

İlk olarak, $\LeftR{\lif}$'in sağ öncülünde biten bir $\RightR{\lforall}$
bulunduğunu varsayalım. Bu durumda $\pi$'nin ilgili alt-!!{proof}{}ı,
$\pi_1$ alt-!!{proof}{}ında şöyle biter:
\[
\AxiomC{}
\RightLabel{$\pi_2$}
\Deduce$\Gamma' \fCenter \Delta', \lforall[x][!A(x)], !B$
\AxiomC{}
\RightLabel{$\pi_3$}
\Deduce$!C, \Gamma' \fCenter \Delta', !A(c)$
\RightLabel{\RightR{\lforall}}
\UnaryInf$!C, \Gamma' \fCenter \Delta', \lforall[x][!A(x)]$
\RightLabel{\LeftR{\lif}}
\BinaryInf$!B \lif !C, \Gamma' \fCenter \Delta', \lforall[x][!A(x)]$
\DisplayProof
\]
$\pi$'nin düzenli olduğunu varsayabiliriz; dolayısıyla özdeğişken~$c$,
$\pi_3$ dışında geçmez; özellikle $!B \lif !C$ içinde ya da $\pi_2$ içinde
geçmez. Şimdi şu $\pi_1'$ ispatını ele alın:
\[
\AxiomC{}
\RightLabel{$\pi_2'$}
\Deduce$\Gamma' \fCenter \Delta', \lforall[x][!A(x)], !A(c), !B$
\AxiomC{}
\RightLabel{$\pi_3$}
\Deduce$!C, \Gamma' \fCenter \Delta', \lforall[x][!A(x)], !A(c)$
\RightLabel{\LeftR{\lif}}
\BinaryInf$!B \lif !C, \Gamma' \fCenter \Delta', \lforall[x][!A(x)], !A(c)$
\RightLabel{\RightR{\lforall}}
\UnaryInf$!B \lif !C, \Gamma' \fCenter \Delta', \lforall[x][!A(x)], \lforall[x][!A(x)]$
\DisplayProof
\]
burada $\pi_2'$, $\pi_2$'nin sağına $!A(c)$ ile zayıflatma uygulanarak elde
edilmiştir. (Bu, hiçbir çıkarım, özellikle dereceyi etkileyebilecek bir
niceleyici çıkarımı eklemez. $!A(c)$ içindeki hiçbir sabit $\pi_2$'de
özdeğişken olarak kullanılmadığından $\pi_2$'deki herhangi bir özdeğişken
koşulunu da bozmaz.) $c$, $!B \lif !C$ içinde geçmediğinden
\RightR{\lforall} üzerindeki özdeğişken koşulu hâlâ sağlanır.

$!B \lif !C, \Gamma' \fCenter \Delta', \lforall[x][!A(x)]$ sekantının bir
$\pi_1''$ !!a{proof}{}ını elde etmek için büzülmenin kabul edilebilirliğine
başvururuz. $\lforall[x][!A(x)]$ için $\RightR{\Contraction}$'ın kabul
edilebilirliğinin ispatı ve bu ispatta kullanılan $\RightR{\lforall}$'ın
tersinirliğinin ispatı hiçbir çıkarımın derecesini değiştirmemiştir: en fazla
bazı çıkarımları silmiştir. Dolayısıyla $\pi_1''$ içinde $\pi_1$'dekinden daha
fazla niceleyici çıkarımı yoktur ve $\pi_1''$ içindeki her niceleyici
çıkarımının altında, $\pi_1$'de ona karşılık gelen çıkarımın altındakiyle aynı
sayıda önermesel çıkarım vardır---son $\RightR{\lforall}$ dışında: $\pi$
içinde altında bir $\LeftR{\lif}$ vardı; bunu $\pi_1''$ içinde onun üstüne
taşıdık. $\pi$ içindeki $\pi_1$'i $\pi_1''$ ile değiştirin. Ortaya çıkan
!!{proof}{}ın derecesi daha düşüktür; çünkü en azından $\RightR{\lforall}$
çıkarımının derecesi (en az)~$1$ azalmıştır. Şimdi tümevarım varsayımını
uygulayın.

Niceleyici çıkarımının $\RightR{\lexists}$ olduğu durumlar daha da kolaydır;
çünkü büzülmeye ihtiyaç yoktur ve hiçbir özdeğişken koşulundan kaygılanmak
gerekmez. Örneğin $\RightR{\lexists}$'in, bu kez sol öncül olarak,
$\RightR{\land}$ tarafından izlendiğini varsayalım. İlgili alt-!!{proof}
$\pi_1$'dir:
\[
\AxiomC{}
\RightLabel{$\pi_2$}
\Deduce$\Gamma' \fCenter \Delta', \lexists[x][!A(x)], !A(t), !B$
\RightLabel{\RightR{\lexists}}
\UnaryInf$\Gamma' \fCenter \Delta', \lexists[x][!A(x)], !B$
\AxiomC{}
\RightLabel{$\pi_3$}
\Deduce$\Gamma' \fCenter \Delta', \lexists[x][!A(x)], !C$
\RightLabel{\RightR{\land}}
\BinaryInf$\Gamma' \fCenter \Delta', \lexists[x][!A(x)], !B \land !C$
\DisplayProof
\]
$\pi$ içindeki $\pi_1$'i şu $\pi_1'$ ile değiştirin:
\[
\AxiomC{}
\RightLabel{$\pi_2$}
\Deduce$\Gamma' \fCenter \Delta', \lexists[x][!A(x)], !A(t), !B$
\AxiomC{}
\RightLabel{$\pi_3'$}
\Deduce$\Gamma' \fCenter \Delta', \lexists[x][!A(x)], !A(t), !C$
\RightLabel{\RightR{\land}}
\BinaryInf$\Gamma' \fCenter \Delta', \lexists[x][!A(x)], !A(t), !B \land !C$
\RightLabel{\RightR{\lexists}}
\UnaryInf$\Gamma' \fCenter \Delta', \lexists[x][!A(x)], !B \land !C$
\DisplayProof
\]
burada $\pi_3'$, $\pi_3$'ün sağına $!A(t)$ ile zayıflatma uygulanarak elde
edilmiştir. Bu zayıflatma yalnızca $\pi_3$ içindeki her sekantın sağına
$!A(t)$ eklemeyi içerir; dolayısıyla hiçbir niceleyici çıkarımının derecesini
değiştirmez. $\pi'$ ile $\pi$ içindeki niceleyici çıkarımlarının dereceleri,
$\RightR{\land}$ ile yer değiştiren ve derecesi~$1$ azalan
$\RightR{\lexists}$ çıkarımı dışında aynıdır. Tümevarım varsayımı uygulanır.
\end{proof}

\begin{prob}
\olref{thm:midsequent} ispatında $\LeftR{\lexists}$'in $\RightR{\lif}$'in
öncülünde bittiği ve $\LeftR{\lforall}$'ın $\LeftR{\lor}$'un sol öncülünde
bittiği durumları betimleyin.
\end{prob}

\begin{proof}[Herbrand Teoremi'nin İspatı]
$\pi$'nin $\Sequent \lexists[x][!A(x)]$ ile bittiğini varsayalım.
\olref{thm:midsequent} gereği $\pi$'yi, şu biçimde olması gereken bir orta
sekant~$S$ içeren bir $\pi'$ !!a{proof}{}ına dönüştürebiliriz:
\[
\Sequent \lexists[x][!A(x)], !A(t_1), \dots, !A(t_n).
\]
$S$'nin üstündeki bütün çıkarımlar önermesel olduğundan
$\lexists[x][!A(x)]$, $S$'nin üstündeki bir çıkarımda ana formül olamaz. Atomik
olmadığından bir aksiyomda da ana formül olamaz. $\lexists[x][!A(x)]$,
$!A(t_1)$'in bir öz alt formülü olamayacağından ($!A(x)$'in niceleyicisiz
olduğunu anımsayın) $S$'nin üstünde etkin de olamaz. Dolayısıyla $S$ ile biten
alt-!!{proof} içinden $\lexists[x][!A(x)]$'in silinmesi, $\Sequent !A(t_1),
\dots, !A(t_n)$'nin doğru bir !!{proof}{}ını verir; buradan $n-1$
$\RightR{\lor}$ çıkarımıyla $\Sequent !A(t_1) \lor \dots \lor !A(t_n)$'nin
bir !!a{proof}{}ı elde edilebilir.
\end{proof}

% \begin{prob}
% Önce $\lexists[x][\lforall[y][(!A(x) \lif !A(y))]]$ için kesmesiz bir
% $\Sequent \lexists[x][\lforall[y][(!A(x) \lif !A(y))]]$ !!{proof}{}ı bulun;
% ardından (gerekirse) bir orta sekant bulmak üzere
% \olref{thm:midsequent}'deki dönüşümü uygulayarak bu formülün bir Herbrand
% ayrışımını bulun.
% \end{prob}

